\chapter*{Glossary}
\addcontentsline{toc}{chapter}{Glossary}

\begin{description}

    \item[monosemanticity] A property of a neuron or feature whose activation corresponds to a single, identifiable concept or pattern, in contrast to polysemantic units that respond to multiple unrelated concepts due to superposition.

  \item[Attribution graph] A directed acyclic graph whose nodes represent model components
    (token embeddings, SAE or CLT features, attention heads, output logits) and whose edges
    represent gradient-weighted information flow between them.
    Introduced as a methodology by \citet{lindsey2025biology}.

  \item[Circuit] A subgraph of the model's computational graph comprising a small set of
    components (attention heads, MLP layers, or individual features) that together implement a
    specific behaviour.
    The term is used in mechanistic interpretability to denote minimal, interpretable
    sub-computations.

  \item[CLT --- Cross-Layer Transcoder] A sparse dictionary-learning module, analogous to a
    sparse autoencoder, that maps MLP inputs at one layer to residual-stream contributions,
    potentially spanning across layers.

  \item[FFN --- Feed-Forward Network] Synonym for MLP in the transformer literature.

  \item[Grokking] The phenomenon whereby a neural network trained on a small dataset first
    memorises the training examples (near-perfect training accuracy, poor generalisation) and
    then, after many additional gradient steps, abruptly transitions to a generalising solution.
    Mechanistically characterised for modular arithmetic by \citet{nanda2023progress}.

  \item[LLM --- Large Language Model] A neural language model with a parameter count
    large enough that emergent capabilities arise from scale alone, typically trained on
    broad natural-language corpora via next-token prediction.

  \item[MLP --- Multilayer Perceptron] In the transformer architecture, the feed-forward
    sub-layer of each block, consisting of two learned linear projections with a pointwise
    nonlinear activation in between.
    Despite the name, transformer MLPs typically have only a single hidden layer.

  \item[RMSNorm --- Root Mean Square Layer Normalisation] A variant of layer normalisation
    that rescales activations by their root mean square without subtracting the mean.
    Used in Qwen and other modern architectures in place of standard LayerNorm.

  \item[SAE --- Sparse Autoencoder] A dictionary-learning module trained to decompose
    neural network activations into a sparse linear combination of learned feature directions.
    Used in mechanistic interpretability to identify human-interpretable features inside a model.

  \item[Lookup table feature] A transcoder feature whose activation matrix
    $M_{\ell,k}[a,b]$ concentrates on the $200$ cells satisfying
    $a \bmod 10 \in \{d_1, d_2\}$ and $b \bmod 10 \in \{d_1, d_2\}$ for a
    specific ones-digit pair $(d_1, d_2)$.
    Such features encode the pre-computed result of the ones-digit sum
    $d_1 + d_2$ and form the mechanistic primitive of addition circuits
    identified by \citet{lindsey2025biology}.

  \item[Operand matrix] A $100 \times 100$ array $M_{\ell,k}[a,b]$ recording the
    activation of transcoder feature $k$ at layer $\ell$ at the final input-token
    position, swept over all integer-pair prompts \texttt{calc:~$a$+$b$=~} with
    $a, b \in \{0, \ldots, 99\}$.
    The geometric pattern of this matrix encodes the quantity that the feature tracks.

  \item[Representation engineering] A methodology for extracting and analysing concept
    directions in the activation space of a neural network by computing the mean difference
    in hidden-state vectors between semantically contrasting inputs.
    Proposed by \citet{zou2023representation} as a top-down approach to model transparency,
    complementary to the bottom-up circuit-based programme.

  \item[Sharpness index] A scalar summary of how tightly a concept delta's norm is concentrated
    around its peak layer, defined as the fraction of total norm mass lying within a window of
    width three centred at the layer of maximum norm.
    Values close to one indicate a localised representation; values close to $3/L$ indicate
    a representation that is uniformly distributed across the network's $L$ layers.


  \item[Activation patching] A causal intervention technique in which a single intermediate
    representation, here the residual-stream vector at a chosen layer and token position,
    is replaced by the corresponding value from a different input.
    The resulting change in output logits quantifies the causal sufficiency of that
    representation for driving the model's prediction.

  \item[Gradient-dot-delta] A first-order approximation to activation patching obtained by
    computing the inner product of the output-margin gradient with respect to a layer's
    residual stream, evaluated at the baseline input, and the concept delta vector at that layer.
    The approximation is exact when the output margin is linear in the residual stream;
    agreement with activation patching in practice confirms that the relevant neighbourhood
    is approximately linear.

  \item[Source-exclusive feature] A transcoder feature that is strongly active under a source
    prompt condition and weakly active under a target condition, qualifying it for suppression
    during a causal intervention scan.  The complementary notion, a target-exclusive feature,
    is one that is strongly active in the target condition and weakly active in the source,
    qualifying it for injection.  Formally, feature $k$ at layer $\ell$ is source-exclusive if
    $a_{\ell,k}^{\,\mathrm{tgt}} < \rho\cdot a_{\ell,k}^{\,\mathrm{src}}$ for a
    threshold $\rho \in (0,1)$.

  \item[Intervention Scan] An experiment in which a causal modification of the model's
    internal activations is applied at a sequence of increasing strengths $\gamma$, and the
    resulting next-token probabilities for a set of tracked tokens are recorded at each value.
    A successful scan produces a monotone crossover between the source-condition and
    target-condition token probability curves, providing evidence that the modified features
    causally mediate the targeted aspect of the output.

  \item[Operand divergence position] The index of the first token at which two prompt
    sequences differ, used in the operand swap experiment to locate the token position at
    which the operand concept enters the context.  For the English prompts \textit{the opposite
    of ``small'' is ``} and \textit{the opposite of ``hot'' is ``}, this is the position of
    the tokens corresponding to ``small'' and ``hot'' respectively.

\end{description}
